\cotefile{16}
\scanpage{53}

Supposons $k\geq0$, et le lemme prouvé pour $k'<k$\unsure{.} On sait déjà que les $p_\ell:H^*(X)\mathbin{-}H^\ell(X)$ sont $\in\operatorname{Alg}$ pour $i<k$, donc (transp.) pour $i>2n-k$. Soit
\[
T_k=\operatorname{id}-\sum_{i<k}p_i-\sum_{j>\unsure{2n-k|2-k}}p_j,
\]
donc $T_k\in A$, plus gén. $T_{k'}\in A$ si $k'\leq k$. Alors $\zeta^{n-k}T_k$ est nul sur $H^i$, $j\neq k$, et induit
\diagram{\begin{tikzcd}[ampersand replacement=\&]
H^k\arrow[r,phantom,"{\unsure{\rightarrow|\twoheadrightarrow}}" description,"{\substack{\zeta^{n-k}\\\sim}}"]\&H^{2n-k}
\end{tikzcd}}
, donc $u_k\zeta^{n-k}T_k=p_k$.

D'autre part, si $i+2j=k$, \inserted{$j>0$,} la projection sur $\zeta^jP^i$ est aussi $\in A$, comme on voit \deleted{\unsure{ainsi \unsure{;|:}|aussitôt}} par récurrence ascendante sur $i$ (donc descendante sur $j$) ainsi : si $i<0$ rien à prouver. Si c'est fait pour les $i'<i$, on forme $\zeta^{n-k+j}T_k$, \deleted{$\unsure{H|\text{il}|\zeta}$} qui induit \unsure{zéro|0} sur $H^\alpha$ pour $\alpha\neq k$, et induit
\diagram{\begin{tikzcd}[ampersand replacement=\&]
H^k\unsure{\ell}\arrow[r,dash,"\zeta^{n-k-j}"]\&H^{2n-k+2j}
\end{tikzcd}}
, \inserted{\unsure{proj|qui} est nul} \deleted{\unsure{sur},} sur les facteurs $\zeta^{j'}P^{i'}$ pour $j'<j$, \inserted{\unsure{de noyau|la somme|des mêmes} $N$,} et sur les facteurs $\zeta^{j'}P^{i'}$ avec $j'\geq j$ (i.e. $i'\leq i$), \deleted{coïncide \unsure{à|est} $1$} dont la somme n'est autre que $\zeta^jH^{k-2j}$, son composé avec l'isom
\diagram{\begin{tikzcd}[ampersand replacement=\&]
H^{k-2j}\arrow[r,"{\substack{\zeta^j\\\unsure{\sim}}}"]\&\zeta^jH^{k-2j}
\end{tikzcd}}
est $\zeta^{n-(k-2j)}$. Donc l'hom $u_{k-2j}\zeta^{\unsure{n-k+j|n-k-j}}T_k$ \unsure{est|$c$|$\epsilon$}, nul sur $N$, et sur $\zeta^jH^{k-2j}$ l'inverse de
\diagram{\begin{tikzcd}[ampersand replacement=\&]
H^{k-2j}\arrow[r,"\sim"]\&\zeta^jH^{k-2j}
\end{tikzcd}}
. Donc

\mspage{52}
